\begin{frame}
\frametitle{Frozen-spin}
\framesubtitle{}

\vspace{-0.3cm}
\begin{block}{Spin-precession of particle MDM   \textit{relative} to direction of flight:}
	\begin{equation}
	\begin{split}
	\vec \Omega & = \vec \Omega_\text{MDM} - \vec \Omega_\text{cyc} \\
	& = - \frac{q}{\gamma m} \left[G \gamma \vec B_\perp + (1 + G) \vec B_\parallel - \left(G\gamma - \frac{\gamma}{\gamma^2 -1} \right) \frac{\vec \beta \times \vec E}{c}\right].  
	\end{split}
	\end{equation}
	\vspace{-0.3cm}
	\begin{itemize}
		\item[$\Rightarrow$] \textbf{$\vec \Omega = 0$ called \textit{frozen spin}, because  momentum and spin stay aligned.}
	\end{itemize}
	\begin{itemize}
		\item In the absence of magnetic fields ($B_\perp = \vec B_\parallel = 0$), 
		\begin{equation}
		\vec \Omega=0, \text{ if } \left(G\gamma - \frac{\gamma}{\gamma^2 -1} \right) = 0.
		\label{eq:magic-momentum}
		\end{equation}
		\vspace{-0.3cm}
		\begin{itemize}
			\item[-] Possible for particles with $G>0$:   proton ($G=1.793$) or electron ($G=0.001$).
		\end{itemize}
	\end{itemize}
\end{block}


\begin{exampleblock}{For protons:  (\ref{eq:magic-momentum}) $\Rightarrow$ \textit{magic momentum}:}
	\begin{equation}
	G- \frac{1}{\gamma^2-1} = 0 \Leftrightarrow G = \frac{m^2}{p^2} 
	\quad \Rightarrow \quad \boxed{p = \frac{m}{\sqrt{G}} = \SI{700.740}{MeV.c^{-1}}}
	\end{equation}
\end{exampleblock}

\end{frame}